\chapter{สมดุลกลและโมเมนต์ของแรง}
\label{ch:ch06_equilibrium}
\index{สมดุลกลและโมเมนต์}

\begin{conceptbox}[title=วัตถุประสงค์การเรียนรู้ประจำบท]
เมื่อสิ้นสุดการศึกษาบทเรียนนี้ นักศึกษาสามารถ
\begin{enumerate}
    \item เมื่อศึกษาบทนี้จบแล้ว ผู้เรียนควรมีความสามารถดังนี้
    \item อธิบายเงื่อนไขสมดุลของอนุภาค (particle equilibrium) และของวัตถุแข็งเกร็ง (rigid body equilibrium) ได้อย่างถูกต้อง
    \item คำนวณโมเมนต์ของแรง (ทอร์ก) รอบจุดหมุนที่กำหนดให้ พร้อมระบุทิศทางของทอร์กด้วยกฎมือขวา
    \item เขียนแผนภาพวัตถุอิสระ (free-body diagram) และตั้งสมการสมดุลแรงและสมดุลโมเมนต์ได้อย่างเป็นระบบ
    \item แก้โจทย์ปัญหาสมดุลของโครงสร้างอย่างง่าย เช่น คานที่รับน้ำหนัก บันไดพิงผนัง และวัตถุแขวนด้วยเชือกหลายเส้น
\end{enumerate}
\end{conceptbox}

\section{ทฤษฎีและหลักการพื้นฐาน}

\begin{bioappbox}
\textbf{ชีวกลศาสตร์ของระบบโครงร่างและคานงัดในสิ่งมีชีวิต}

โครงสร้างกระดูก ข้อต่อ และกล้ามเนื้อของมนุษย์และสัตว์ทำหน้าที่เป็นระบบคานงัดทางกลศาสตร์ (Biomechanics Levers) เช่น การยกแขนโดยอาศัยแรงดึงของกล้ามเนื้อไบเซปส์ (Biceps) ซึ่งเป็นคานประเภทที่สามที่เสียเปรียบเชิงกลแต่ได้เปรียบด้านความเร็วและระยะการเคลื่อนที่ การคำนวณจุดศูนย์ถ่วง (Center of Gravity) และโมเมนต์ของแรง $\sum \tau = 0$ ช่วยอธิบายเสถียรภาพในการทรงตัวของสัตว์และการออกแบบสรีรศาสตร์ในการทำงานเกษตรกรรมเพื่อลดการบาดเจ็บของกล้ามเนื้อ
\end{bioappbox}

\begin{labbox}
1. ตรวจสอบตำแหน่งศูนย์ (Zero Error) ของเครื่องมือวัดทุกชนิดก่อนเริ่มบันทึกข้อมูล
2. ทำการวัดซ้ำอย่างน้อย 3 ถึง 5 ครั้งในแต่ละจุดทดลองเพื่อคำนวณหาค่าเฉลี่ยและค่าเบี่ยงเบนมาตรฐาน
3. บันทึกผลการวัดพร้อมระบุหน่วยเอสไอ (SI Units) และค่าความคลาดเคลื่อนของการวัดเสมอ
\end{labbox}

\section{ตัวอย่างโจทย์และการวิเคราะห์เชิงฟิสิกส์}

\begin{examplebox}
โจทย์เกี่ยวกับคาน (beam) และบันได (ladder) เป็นตัวแทนคลาสสิกของปัญหาสมดุลวัตถุแข็งเกร็งที่พบได้บ่อยที่สุด และเป็นตัวอย่างที่ดีที่แสดงให้เห็นถึงประโยชน์ของการเลือกจุดหมุนอย่างชาญฉลาด

สำหรับปัญหาคานที่วางพาดอยู่บนฐานรองรับสองจุดและรับน้ำหนักบรรทุกหลายตำแหน่ง หากต้องการหาแรงปฏิกิริยาที่จุดรองรับจุดใดจุดหนึ่งโดยไม่ต้องแก้ระบบสมการพร้อมกัน กลยุทธ์ที่มีประสิทธิภาพคือการเลือกจุดหมุนให้ตรงกับตำแหน่งของแรงปฏิกิริยาอีกจุดหนึ่งที่ยังไม่ทราบค่า เนื่องจากแรงปฏิกิริยาที่จุดหมุนนั้นจะมีแขนโมเมนต์เป็นศูนย์ จึงไม่ปรากฏในสมการโมเมนต์ ทำให้สามารถแก้หาแรงปฏิกิริยาอีกจุดหนึ่งได้โดยตรงจากสมการ $\Sigma$$\tau$ = 0 เพียงสมการเดียว จากนั้นจึงนำค่าที่ได้ไปแทนในสมการ $\Sigma$F = 0 เพื่อหาแรงปฏิกิริยาที่จุดที่เหลือ

สำหรับปัญหาบันไดที่พิงผนัง จุดหมุนที่สะดวกที่สุดมักเป็นจุดที่เท้าบันไดสัมผัสพื้น เพราะที่จุดนี้มีแรงปฏิกิริยาจากพื้นและแรงเสียดทานที่พื้นกระทำอยู่ ซึ่งทั้งสองแรงนี้ผ่านจุดหมุนพอดี จึงมีแขนโมเมนต์เป็นศูนย์และไม่ปรากฏในสมการโมเมนต์ ทำให้สามารถหาแรงปฏิกิริยาจากผนังได้โดยตรง ก่อนนำไปหาแรงเสียดทานที่พื้นจากสมการ $\Sigma$Fx = 0 ต่อไป ข้อควรระวังสำคัญคือ หากผนังลื่น (ไม่มีแรงเสียดทาน) แรงปฏิกิริยาจากผนังจะมีเฉพาะองค์ประกอบแนวราบเท่านั้น ในขณะที่พื้นซึ่งมีความฝืดจะมีทั้งแรงปฏิกิริยาตั้งฉาก (แนวดิ่ง) และแรงเสียดทาน (แนวราบ) ร่วมกัน

หลักการเลือกจุดหมุนอย่างเหมาะสมนี้ไม่เพียงลดความซับซ้อนทางคณิตศาสตร์ แต่ยังช่วยลดโอกาสเกิดความผิดพลาดในการคำนวณ และเป็นทักษะสำคัญที่ผู้เรียนควรฝึกฝนให้เชี่ยวชาญก่อนที่จะนำไปประยุกต์ใช้กับโครงสร้างที่ซับซ้อนยิ่งขึ้นในวิชาต่อ ๆ ไป
\end{examplebox}

\section{แบบฝึกหัดทบทวนและคำถามท้ายบท}

\begin{enumerate}
    \item จงอธิบายความแตกต่างระหว่างเงื่อนไขสมดุลของอนุภาคกับเงื่อนไขสมดุลของวัตถุแข็งเกร็ง
    \item วัตถุมวล 8 กิโลกรัม แขวนนิ่งด้วยเชือกสองเส้นที่ทำมุม 45$^\circ$ เท่ากันทั้งสองข้างกับแนวราบ จงหาแรงตึงในเชือกแต่ละเส้น (g = 9.8 m/s$^2$)
    \item แรงขนาด 60 นิวตัน กระทำที่ปลายประแจ ห่างจากจุดหมุนน็อต 0.25 เมตร โดยทำมุม 90$^\circ$ กับด้ามประแจ จงหาขนาดของทอร์กที่เกิดขึ้น
    \item หากแรงในข้อ 3 กระทำทำมุม 30$^\circ$ กับด้ามประแจแทนที่จะเป็น 90$^\circ$ ทอร์กที่ได้จะมีค่าเท่าใด และมีค่าน้อยกว่าเดิมกี่เท่า
    \item คานสม่ำเสมอยาว 8 เมตร หนัก 400 นิวตัน วางบนฐานรองรับที่ปลายทั้งสองข้าง มีน้ำหนัก 600 นิวตัน วางห่างจากปลายซ้าย 3 เมตร จงหาแรงปฏิกิริยาที่ฐานรองรับทั้งสองจุด
    \item บันไดยาว 6 เมตร หนัก 200 นิวตัน พิงผนังลื่นทำมุม 65$^\circ$ กับพื้น จงหาแรงปฏิกิริยาจากผนังและแรงเสียดทานที่พื้น
\end{enumerate}

% สิ้นสุดบทเรียน
